\section{Scope, conventions, and assumptions}
\label{sec:scope}

This section is the contract for everything that follows. SI units
throughout; $\rho_0=\SI{1.20}{\kilogram\per\cubic\meter}$,
$c=\SI{343}{\meter\per\second}$. Sound pressure levels are dB re
$p_0=\SI{20}{\micro\pascal}$ computed from \emph{rms} pressure, half-space
radiation, on axis.

\begin{center}
	\small
	\begin{tabular}{@{}llll@{}}
	\toprule
		$F_s$ & free-air resonance & $\Mms$ & moving mass\\
		$\Qes,\Qms,\Qts$ & Q factors & $\Re$ & DC resistance\\
		$\Vas$ & compliance vol. & $\Bl$ & force factor\\
		$\Sd$ & piston area & $\Le$ & coil inductance\\
		$\Xmax$ & linear excursion & $\Pmax$ & rated power\\
		$\Vd$ & displacement vol. & $\Vb$ & box volume\\
		$F_c,\Qtc$ & in-box corner, $Q$ & $f_L$ & inductive screen\\
	\bottomrule
	\end{tabular}
\end{center}

\subsection{Reference bases}
\label{sec:scope-bases}

Four distinct pressure references are used, and no figure or table mixes
them:

\begin{enumerate}
\item[B1.] \emph{Intrinsic ceiling}: one driver, half space,
	$r=\SI{1}{\meter}$. Used only to characterize a driver, never to assess
	a system.
\item[B2.] \emph{Couch-area reference demand}: the level required at the
	declared listening distance $r_{\mathrm{listen}}$ over the primary
	listening area. All feasibility checks use this basis.
\item[B3.] \emph{Mono manifold total}: the aggregate level of the complete
	low-bass manifold at B2. The per-driver share of the required displaced
	volume is the total divided by the number of manifold drivers.
\item[B4.] \emph{One independent stereo channel}: the level of one
	upper-bass channel at B2, with no credit from the opposite channel.
\end{enumerate}

Electrical quantities are per \emph{physical} driver, since each has its
own amplifier channel. Peak and rms conventions are stated at each
equation: displacements and the displaced volume $V_{\mathrm{dem}}$ are
peak values, while voltage, current, and power are rms over the evaluated
waveform.

\subsection{Canonical scenario}
\label{sec:scope-scenario}

The worked scenario used for every figure and table is:

\begin{center}
	\small
	\begin{tabular}{@{}p{2.5cm}p{5.0cm}@{}}
	\toprule
		room & $V_{\mathrm{room}}=\SI{60}{\cubic\meter}$,
			$L_{\max}=\SI{6}{\meter}$\\
		room model & leaky pressure zone,
			$f_\ell=\SI{10}{\hertz}$\\
		listening distance & $r_{\mathrm{listen}}=\SI{3}{\meter}$\\
		bands & sub $15$--$\SI{80}{\hertz}$;
			upper bass $80$--$\SI{250}{\hertz}$\\
		targets & \SI{110}{dB} manifold total (B3);\\
		 & \SI{105}{dB} per stereo channel (B4)\\
		amplifier, per driver & $\SI{90}{\volt}$ rms, $\SI{15}{\ampere}$ rms,\\
		 & \SI{500}{\watt} continuous, \SI{1}{\kilo\watt} burst\\
		alignment ceiling & $\Qtc\le0.55$\\
		enclosure limit & $\Vb\le\min(10\Vas,\ \SI{1.0}{\cubic\meter}/N)$
			per driver; \SI{1.0}{\cubic\meter} total per role\\
		excursion & $\Xmax$ hard limit;
			$0.80$ preferred margin\\
		Doppler screen & first-sideband ratio $\le0.02$\\
		transient & one-cycle rectangular burst, 8 phases\\
	\bottomrule
	\end{tabular}
\end{center}

Target levels are maximum-output design tests. The pressure-zone corner of
this room is $f_{\mathrm{pz}}=c/2L_{\max}\simeq\SI{28.6}{\hertz}$.

\subsection{Assumptions}
\label{sec:scope-assumptions}

\begin{enumerate}
\item[A1.] \emph{Lumped-parameter LTI small-signal model.} The screening
	domain ends at $|x|=\Xmax$; operation below that boundary does not by
	itself establish linear behavior. Cone breakup is a $ka\gtrsim1$
	phenomenon above the bands considered here, so the rigid-piston premise
	is the same one the published T/S parameters already presume.

\item[A2.] \emph{Coil inductance neglected in the transfer relations.}
	$f_L=\Re/(2\pi\Le)$ is retained as a screening frequency only. $\Le$ is
	itself frequency- and level-dependent, so $f_L$ is a rough indicator of
	where the model degrades, not a sharp validity boundary.

\item[A3.] \emph{Baffled rigid piston, far field, $ka\lesssim1$}
	($a=\sqrt{\Sd/\pi}$): the monopole approximation underlying every SPL
	relation below.

\item[A4.] \emph{Lossless sealed box}, adiabatic air spring. Box losses are
	a second-order correction that the ranking does not depend on.

\item[A5.] \emph{$\Pmax$ is a continuous coil-dissipation rating}
	$\langle i^2\rangle\Re$. Published ratings follow different standards
	\cite{aes2}, so comparisons are meaningful only within one standard, and
	a rating uncertainty of a few dB is propagated explicitly
	(\cref{sec:procedure}). Thermal compression \cite{button1992} raises
	$\Re$ during sustained drive and detunes $\Qtc$ and $F_c$; it is
	reported as a risk indicator, not modelled.

\item[A6.] \emph{Amplifier limits are explicit inputs.} Voltage, current,
	continuous power, and burst power per physical driver are hard gates
	(\cref{eq:ampgates}). The scenario is considered viable only if the
	driver, enclosure, and room constraints bind before those limits.

\item[A7.] \emph{$\Xmax$ is used only as the clipping boundary.} A
	distortion-defined value \cite{iec62458,klippel2003xmax} is preferred;
	geometric values must be converted or flagged. Excursion utilization
	$\xi_x=X_{\mathrm{dem}}/\Xmax$ is a risk proxy, not predicted THD.
	Public large-signal measurements supersede the proxy wherever they
	exist.

\item[A8.] \emph{Room effects are separated from free-field ceilings.} The
	ceilings of \cref{sec:freq} are free-field driver properties; room
	demand enters only through \cref{sec:room}.

\item[A9.] \emph{Summation is restricted to the mono manifold.} The
	manifold's drivers carry one common signal in a compact arrangement
	(spacing $\ll$ wavelength, consistent with A3), so their volume
	velocities are taken as coherent at the couch-area reference: each
	driver supplies $1/N$ of the required displaced volume. Independent
	stereo channels receive no summing credit. Total radiated power is
	\emph{not} inferred from on-axis pressure summation, mutual radiation
	impedance is not modelled, and the electrical demand of each physical
	driver is computed from its own displacement across the complete role
	band.

\item[A10.] \emph{Nonlinear indicators are proxies.} They are
	leading-order, dimensionless screening quantities reported separately.
	They are not summed, and they are not converted into an absolute
	distortion percentage.

\item[A11.] \emph{DSP correction is finite-band and regularized.} Linear
	correction trades response error against displacement, voltage, current,
	thermal load, noise, and model error; it does not reduce nonlinear
	residuals, and its practical limits over an area are those reported in
	\cite{makivirta2003,cecchi2018}.
\end{enumerate}

\subsection{Limitations of public data}
\label{sec:scope-datalimits}

Datasheet fields carry mixed conventions for $\Xmax$, $\Pmax$, and $\Le$,
occasional missing fields, and no statement of unit-to-unit spread. The
procedure therefore (i) rejects records missing a required field rather
than imputing it, (ii) records which fields were derived rather than
published, and (iii) reports rank stability under declared perturbations of
$\Xmax$ and $\Pmax$ instead of presenting a single ranking as exact.
